% =============spacing=============
% reduce space 
% \renewcommand{\baselinestretch}{0.9745}

% \newcommand{\controlspace}{
% \newcommand{\authsep}{\hspace{1em}}
% \newcommand{\instsep}{\hspace{4em}}
% \newcommand{\emailsep}{\hspace{2em}}

% \captionsetup[figure]{font=small}
% \captionsetup[table]{font=footnotesize}


% =============Definition=============

% \DeclareMathOperator{\argmin}{argmin}

% \newcommand{\overbar}[1]{\mkern 1.5mu\overline{\mkern-1.5mu#1\mkern-1.5mu}\mkern 1.5mu}

\newcommand{\red}[1]{{#1}}
\newcommand{\todo}[1]{\textcolor{red}{$\Rightarrow$ TODO: #1}}
% \newcommand{\red}[1]{\textcolor{red}{#1}}

% Small Title
% \newcommand{\stitle}[1]{\vspace{0.1ex} \noindent {\bf {#1}}}
% \newcommand{\sstitle}[1]{\vspace{0.1ex} \noindent{\textit {\underline {#1}}}}
% \newcommand{\ssstitle}[1]{\vspace{0.1ex} \noindent{\textbf {#1}}}

\newcommand{\stitle}[1]{{\noindent \bf {#1}}}
\newcommand{\sstitle}[1]{{\noindent \textit {\underline {#1}}}}
\newcommand{\ssstitle}[1]{{\noindent \textit {#1}}}

\newcommand{\kw}[1]{{\ensuremath {\mathsf{#1}}}\xspace}
\newcommand{\bkw}[1]{{\ensuremath {\mathsf{\textbf{#1}}}}\xspace}
\newcommand{\bkwnospace}[1]{{\ensuremath {\mathsf{\textbf{#1}}}}}
\newcommand{\kwnospace}[1]{{\ensuremath {\mathsf{#1}}}}

\newcommand{\pluseq}{\mathrel{+}=}
\newcommand{\asteq}{\mathrel{*}=}

\newcommand{\tree}{\mathcal{B}}

\newcommand{\cpi}{\kw{CPI}}
\newcommand{\sct}{\kw{SCT}}
% \newcommand{\tree}{\mathcal{CPI}}
\makeatletter
% 定义主命令：看下一个字符是不是“{”
\def\TPath{%
  \@ifnextchar\bgroup        % 如果下一个是 { → 执行 \TPathWithArg
    {\TPathWithArg}%
    {\TPathNoArg}%           % 否则执行 \TPathNoArg
}

% 带参数的分支
\def\TPathWithArg#1{%
  \mathcal{P}_{#1}%
}

% 不带参数的分支
\def\TPathNoArg{%
  \mathcal{P}%
}
\makeatother

\newcommand{\kcore}{$k$-Core\xspace}
\newcommand{\hyperCD}{\hyperref[alg:HyperCD]{\kwnospace{HyperCD}}\xspace}
\newcommand{\hyperCDM}{\hyperref[alg:CoreEvaluationFinal]{\kwnospace{HyperCD^{+}}}\xspace}
\newcommand{\hyperCDP}{\hyperref[alg:HyperCD+]{\kwnospace{HyperCD^{++}}}\xspace}
\newcommand{\hyperCDF}{\hyperref[alg:HyperCDFinal]{\kwnospace{HyperCD^*}}\xspace}
\newcommand{\hyperCDPar}{\hyperref[alg:HyperCDPar]{\kwnospace{HyperCDPar}}\xspace}

\newcommand{\bhyperCD}{\hyperref[alg:HyperCD]{\bkwnospace{HyperCD}}\xspace}
\newcommand{\bhyperCDP}{\hyperref[alg:HyperCD+]{\bkwnospace{HyperCD^{++}}}\xspace}
\newcommand{\bhyperCDF}{\hyperref[alg:HyperCDFinal]{\bkwnospace{HyperCD^*}}\xspace}
\newcommand{\bhyperCDPar}{\hyperref[alg:HyperCDPar]{\bkwnospace{HyperCDPar}}\xspace}

\newcommand{\maintainSup}{\hyperref[alg:MaintainSup]{\kwnospace{MaintainSup}}\xspace}
\newcommand{\coreCorrection}{\hyperref[alg:CoreEvaluationP]{\kwnospace{CoreCorrection}}\xspace}
\newcommand{\coreEvaluation}{\hyperref[alg:CoreEvaluationP]{\kwnospace{CoreEvaluation}}\xspace}
\newcommand{\coreEvaluationP}{\hyperref[alg:CoreEvaluationFinal]{\kwnospace{CoreEvaluation^*}}\xspace}
\newcommand{\bcoreCorrectionF}{\hyperref[alg:CoreEvaluationFinal]{\bkwnospace{CoreEvaluation^*}}\xspace}

\newcommand{\localcore}{\kw{LocalCore}}
\newcommand{\blocalcore}{\bkw{LocalCore}}
\newcommand{\localcoreo}{\kw{LocalCore}-\kw{OPT}}
\newcommand{\blocalcoreo}{\bkw{LocalCore}-\kw{OPT}}
\newcommand{\nbrkcore}{\kw{nbr}-$k$-\kw{Core}}
\newcommand{\nbrkcores}{\kw{nbr}-$k$-\kw{Cores}}
\newcommand{\bnbrkcore}{\bkw{nbr}-$k$-\bkw{Core}}
\newcommand{\nbrkcorek}[1]{\kw{nbr}-${#1}$-\kw{Core}}

% cnd
\newcommand{\cbnd}{\kw{CND}}
\newcommand{\arb}{\kw{ARB}}
\newcommand{\nucleus}{\kw{ND}}
\newcommand{\nucleusP}{\kw{PND-peel}}
\newcommand{\nucleusOPT}{\kw{PND-Opt}}

% --- relation symbols for path/clique relations ---
\providecommand{\encodedby}{\mathrel{\trianglelefteq}} % strong: "encoded by" (must include holds)
\providecommand{\coveredby}{\mathrel{\sqsubseteq}}     % weak: "covered on the path" (vertex-set inclusion)
% \newcommand{\hold}{\textsf{\textbf{hold}}\xspace}
% \newcommand{\pivot}{\textsf{\textbf{pivot}}\xspace}

\newcommand{\CW}{\hyperref[def:coreW]{\phi}}
\newcommand{\CWs}{\hyperref[def:cis]{\Phi}}
% \ref{the:LB}
\newcommand{\LB}{\hyperref[the:LB]{\kw{LB}}\xspace}


\newcommand{\coreW}{\textit{Core Weight}\xspace}
% \newcommand{\coreW}{\texorpdfstring{\textit{Core Weight}\hypertarget{def:coreW}{}}{Core Weight}\xspace}
\newcommand{\coreWS}{\textit{Core Weight Set}\xspace}
% \ref{def:coreW}

% datasets
\newcommand{\amazon}{\kw{Amazon}}
\newcommand{\dblpFull}{\kw{coauth-DBLP-FULL}}
% \newcommand{\threadsStackOverflow}{\kw{hreads-stack-overflow}}
% % tags-stack-overflow
% \newcommand{\tagsStackOverflow}{\kw{tags-stack-overflow}}
\newcommand{\daEN}{\kw{EN}}
\newcommand{\daHI}{\kw{HI}}
\newcommand{\daCO}{\kw{CO}}
\newcommand{\daGE}{\kw{GE}}
\newcommand{\daDB}{\kw{DB}}
\newcommand{\daDF}{\kw{DF}}
\newcommand{\daTHS}{\kw{THS}}
\newcommand{\daAM}{\kw{AM}}
\newcommand{\daTAS}{\kw{TAS}}
\newcommand{\daAMZ}{\kw{AMZ}}
\newcommand{\daSOA}{\kw{SOA}}
\newcommand{\daPOWS}{\kw{POW}-\kw{S}}
\newcommand{\daPOWL}{\kw{POW}-\kw{L}}
\newcommand{\RUC}{\hyperref[the:onlyChange]{\kw{RUC}}\xspace}
\newcommand{\SUC}{\hyperref[the:sup]{\kw{SUC}}\xspace}
% Theorem
% \newtheorem{theorem}{Theorem}[section]
% \newtheorem{lemma}[theorem]{Lemma}

% \newtheorem{property}{Property}[section]

% \newtheorem{theorem}{Theorem}[section]
% \newtheorem{lemma}[theorem]{Lemma}
% \newtheorem{property}[theorem]{Property}

% \theoremstyle{definition}
% \newtheorem{definition}{Definition}[section]

% % \theoremstyle{definition}
% \theoremstyle{remark}
% \newtheorem{example}{Example}

% \theoremstyle{remark}
% \newtheorem{observation}{Observation}[section]

% \theoremstyle{remark}
% \newtheorem{gen_rule}{Generation Rule}[section]

% \newtheorem*{remark}{Remark}
% \newenvironment{remark}{\begin{itshape}}{\end{itshape}}
% \newenvironment{remark}{\begin{itshape}}{\end{itshape}}
% \newenvironment{remark}{\par\noindent\stitle{Remark.} \itshape}{\par\smallskip}
\NewDocumentEnvironment{remark}{o}
    {\par\noindent\textit{Remark}\IfValueT{#1}{ \textit{(#1)}}. \itshape}
  {\par\smallskip}
% \theoremstyle{definition}
% \newtheorem*{remark}{Remark}


% === Theorem-like environments (independent counters; reset per section) ===
\theoremstyle{plain}
\newtheorem{theorem}{Theorem}[section]
\theoremstyle{plain}
\newtheorem{lemma}{Lemma}[section]
\theoremstyle{plain}
\newtheorem{property}{Property}[section]
\theoremstyle{plain}
\newtheorem{corollary}{Corollary}[section]
\theoremstyle{definition}
\newtheorem{definition}{Definition}[section]
\theoremstyle{remark}
\newtheorem{example}{Example}[section]


\newcommand{\cfig}{Fig.~}
\newcommand{\ctab}{Table~}
\newcommand{\csec}{Section~}
\newcommand{\cdef}{Definition~}
\newcommand{\cthm}{Theorem~}
\newcommand{\clem}{Lemma~}
\newcommand{\cequ}[1]{Equation~(#1)}

% \long\def\comment#1{}


% =============defined for reference======
\newcommand{\reffig}[1]{Fig.~\ref{fig:#1}}
\newcommand{\refsec}[1]{Section~\ref{sec:#1}}
\newcommand{\reftable}[1]{Table~\ref{tab:#1}}
\newcommand{\refalg}[1]{Algorithm~\ref{alg:#1}}
\newcommand{\refeq}[1]{Equation~\ref{eq:#1}}
\newcommand{\refdef}[1]{Definiton~\ref{def:#1}}
\newcommand{\refthm}[1]{Theorem~\ref{thm:#1}}
\newcommand{\reflem}[1]{Lemma~\ref{lem:#1}}
\newcommand{\refrem}[1]{Remark~\ref{rem:#1}}
\newcommand{\refcoro}[1]{Corollary~\ref{coro:#1}}
\newcommand{\refex}[1]{Example~\ref{ex:#1}}
\newcommand{\refppt}[1]{Property~\ref{ppt:#1}}


% Do not define a theorem-style here; a custom Remark environment is provided below.

% =============hold and pivot styling (sans-serif, not bold)=============
\newcommand{\hold}{\textsf{hold}\xspace}
\newcommand{\pivot}{\textsf{pivot}\xspace}
% =============hold and pivot styling (sans-serif, not bold)=============
\newcommand{\holds}{\textsf{holds}\xspace}
\newcommand{\pivots}{\textsf{pivots}\xspace}


% =============defined for Exp======
% 定义主计数器
\newcounter{expcounter}
\setcounter{expcounter}{0}

% 定义子计数器
\newcounter{expsubcounter}[expcounter] % 每次expsection增加时自动重置子计数器
\setcounter{expsubcounter}{0}

% 定义Exp Section
\newcommand{\expsection}[1]{%
    \stepcounter{expcounter}% 增加主计数器
    \setcounter{expsubcounter}{0}% 每次新expsection重置子计数器
    \stitle{Exp-\theexpcounter: #1.}%
}

% 定义Exp Subsection
\newcommand{\expsubsection}[1]{%
    \stepcounter{expsubcounter}% 增加子计数器
    \ssstitle{Exp-\theexpcounter.\theexpsubcounter: #1.}%
}

% Define the \myremove command to strike through text, handling math mode
\newcommand{\myremove}[1]{
  \ifmmode
    % \textcolor{red}{\cancel{\text{$#1$}}}
  \else
    % \textcolor{red}{\sout{#1}}
  \fi
}
% \newcommand{\myremove}[1]{\textcolor{red}{\cancel{#1}}}

% Define the \myadd command to underline added text, handling math mode
\usepackage{xcolor}
\usepackage[normalem]{ulem}

\newcommand{\myadd}[1]{%
  \ifmmode%
    % 数学模式：保持下划线和颜色，不支持换行（公式一般也不换段）
    \textcolor{blue}{\underline{#1}}%
  \else%
    % 文本模式：使用开关式写法 {\color{blue} ...}
    % 这种写法支持空行分段，不会报错
    {\color{blue}#1}%
  \fi%
}
\newcommand{\Rtag}[1]{\textbf{\textcolor{black}{[R-#1]}}~}


% \newcommand{\myaddRA}[1]{%
%   \ifmmode%
%     % 数学模式：保持下划线和颜色，不支持换行（公式一般也不换段）
%     \textcolor{red}{\underline{#1}}%
%   \else%
%     % 文本模式：使用开关式写法 {\color{purple} ...}
%     % 这种写法支持空行分段，不会报错
%     {\color{red}#1}%
%   \fi%
% }

% Revision highlight helper: handles math/text coloring, captions, and table rules
\newcommand{\myaddRBase}[2]{%
  \begingroup
  \ifmmode
    % \underline{#2}%
    #2
  \else
    % \color{#1}% body text in group
    % \arrayrulecolor{#1}% table lines in group
    % \captionsetup{labelfont={color=#1},textfont={color=#1}}%
    % \captionsetup[subfigure]{labelfont={color=#1},textfont={color=#1}}%
    #2%
  \fi
  \endgroup
}

% Revision highlights A/B/C
\newcommand{\myaddRA}[1]{\myaddRBase{red}{#1}}
\newcommand{\myaddRB}[1]{\myaddRBase{orange}{#1}}
\newcommand{\myaddRC}[1]{\myaddRBase{blue}{#1}}



% Define the \myupdate command for old and new text, handling both text and math mode
\newcommand{\myupdate}[2]{%
  \ifmmode
    % \textcolor{red}{\cancel{\text{$#1$}}} 
    % \textcolor{blue}{#2}
    #2
  \else
    % \textcolor{red}{\sout{#1}} 
    % \textcolor{blue}{}
    #2
  \fi
}
